\section{Introduction}
%Graphs represent relationships among entities as vertices connected by edges. Graph algorithms are computational techniques for analyzing these connections and extracting both structural and semantic insights.
% Definition of Traversed-based graph analytics: visit more and more vertices by always operating on a subset of vertices
% Traversal-based graph analytics are a fundamental graph workload for analyzing connectivity and graph structure, and they appear in many application domains. 
\subsection{Traversal-based Graph Analytics}
\label{subsec:intro-traversal}
Traversal-based graph algorithms form a fundamental family of graph analytics methods. 
% They start from one or more source vertices and incrementally discover unvisited vertices by traversing edges, stopping when a specified termination condition is met.
They begin from one or more source vertices as an initial frontier and progressively explore the graph by expanding this frontier with more vertices until a specified termination condition is met.
% Application of traversed-based graph analytics
Representative examples include BFS, SSSP, and MSSP, which are widely used in network analysis to compute reachability and shortest-path distances from one or multiple sources in communication and infrastructure networks \cite{hoefler2009optimized, guimera2005worldwide, zhan1998shortest}. 
In social-network analysis, traversal-based methods support influence estimation, central-actor identification, and diffusion modeling \cite{basuchowdhuri2019fast, pv2024degree, byun2019chronograph}. 
They are also core primitives in AI and knowledge-graph reasoning, enabling state-space search, symbolic inference, and path-based reasoning \cite{russell1995modern, osipjan2025graphtrace, canesche2020traversal, jia2025neural}. 
In scientific computing and bioinformatics, traversal-based algorithms are used to analyze interaction graphs (e.g., molecular and protein–protein interaction networks) and biological pathways \cite{alon2007network, rubel2020augmenting, franzese2019hypergraph}.

\subsection{The Blossom Algorithm and its Parallel Processing}
% Definitions of Maximum Matching .
Maximum matching is a fundamental problem in graph theory and a cornerstone of combinatorial optimization \cite{plummer1986matching}.
It has broad real-world applications, including logistics planning for vehicle-to-request assignment \cite{meshkani2022generalized, abeywickrama2021optimizing, kalantari2023trailer}, large-scale clustering and data analysis such as representative structure construction, entity alignment, and community detection \cite{auer2012graph, dhillon2007weighted, zhu2024survey, sun2018bootstrapping, majmudar2020provable, rezaei2016set, gmati2024new}, cloud and edge job scheduling for task–resource allocation \cite{herodotou2021trident, yang2021edge, ye2024deep}, as well as machine learning \cite{he2021learnable, li2022sigma, wang2021joint}, recommendation systems \cite{han2021graph, naini2015you, su2021neural}, and health resource allocation \cite{roth2005pairwise}.

% Edmond's Blossom Algorithm
Given a graph, a matching is a set of pairwise non-adjacent edges and a maximum matching is a matching of maximum cardinality.
Edmonds’ blossom algorithm \cite{edmonds1965maximum, edmonds1965paths} is the first polynomial-time method for computing maximum matchings in general graphs.
It iteratively searches for augmenting paths and contracts odd-length cycles (blossoms) encountered during the search.
% Traversal-based nature of blossom algorithm
Graph primitives can be broadly classified by execution pattern into iterative algorithms and traversal-based algorithms.
Iterative graph algorithms, such as Triangle Counting and PageRank, apply uniform operations to all vertices in the graph until a convergence criterion is satisfied. 
In contrast, traversal-based algorithms follow the frontier-expansion paradigm described in \S~\ref{subsec:intro-traversal}.
Examining the execution pattern of the blossom algorithm, it proceeds in rounds consisting of expansion, augmenting-path detection, and blossom processing.
The expansion phase directly corresponds to the traversal process, during which unvisited vertices are explored and added to the frontier. 
Augmenting-path detection acts as a termination condition for the traversal, signaling when a valid augmenting path has been found. 
Blossom processing serves as an auxiliary operation that identifies and handles odd cycles, allowing the search to ultimately discover an augmenting path. 
Therefore, the blossom algorithm can be characterized as a traversal-based graph primitive, as its execution pattern closely aligns with the notion of graph traversal.

% Sequential Blossom is slow
Edmonds’ blossom algorithm is computationally expensive, with worst-case time complexity $O(V^2E)$, making it a bottleneck for modern large-scale graphs with millions of vertices and edges.
% Moore's Law => Parallel
While extensive research efforts have optimized the sequential implementation of the blossom algorithm \cite{blum2015maximum, gabow1976efficient, gabow1984efficient, gabow1990data, gabow1991faster}, as Moore’s Law approaches its physical limits, the rate of improvement in single-core performance has largely diminished.
Therefore, developing efficient parallel algorithms is increasingly important for large-scale maximum matching.
% \cite{fan2024x}.
% GPUs => Massively Parallel
Moreover, modern hardware accelerators, particularly GPUs, have been widely used to accelerate graph algorithms by leveraging thousands of cores and high memory bandwidth. 
This raises the question of whether such architectures can also be harnessed to speed up large-scale blossom-based maximum matching.


% X-Blossom, X-Blossom-Pro and X-Blossom++
X-Blossom \cite{fan2025x} is a parallel framework for the blossom algorithm that avoids blossom contraction and dynamic graph modifications, enabling concurrent growth of search structures and the discovery of multiple augmenting paths per phase. 
% \cite{fan2025x, wang2019sep, wang2021grus}.
However, the current implementation of X-Blossom has two key limitations: it aggressively discards reusable search structures, and its CPU-based execution prevents it from fully exploiting the framework's inherent parallelism.
To address these limitations, we develop two enhanced realizations of X-Blossom.
First, we introduce a selective-reset mechanism that clears only the search structures involved in discovered augmenting paths while preserving all other reusable components.
This yields an improved CPU implementation, X-Blossom-Pro, which achieves a speedup of up to $3.26\times$ over the original version and provides a stronger baseline.
Building on X-Blossom-Pro, we further propose X-Blossom++, the first GPU-based realization of the blossom algorithm.
X-Blossom++ incorporates an additional restructuring phase to transform the framework’s irregular parallelism into compact, fine-grained tasks suitable for massive GPU execution. 
This unlocks substantially greater parallelism and delivers up to $39.09\times$ acceleration over X-Blossom-Pro across all datasets.

\subsection{Our Profiling Approach for Performance Insights}
Bulk Synchronous Parallel (BSP) models parallel programs as supersteps of local computation, communication, and synchronization \cite{valiant1990bridging}.
Traversal-based graph processing naturally follows this structure, as each iteration updates active vertices or partitions, exchanges frontier or dependency information, and synchronizes before the next iteration.

However, BSP is architecture-independent and does not directly capture the tightly coupled memory systems and execution mechanisms of modern multi-core CPUs and GPUs.
% For example, A single multi-core CPU system can integrate up to more than 100 nodes with a large L3 cache, and a huge shared memory \cite{cpu_arch_amd}. 
% A single GPU card features thousands of CUDA cores and a large capacity of device memory \cite{nvidia_ampere_architecture_whitepaper}. 
For example, modern CPUs provide many complex cores with a large shared \phantomsection\label{para:first-llc}\textcolor{blue}{Last-Level Cache (LLC)} \cite{cpu_arch_amd}, while GPUs provide thousands of lightweight cores and high-bandwidth device memory \cite{nvidia_ampere_architecture_whitepaper}.
% In such environments, inter-worker communication is largely realized within local computation as shared-memory reads and writes, with fine-grained atomic operations enforcing synchronization. As a result, communication overhead no longer needs to be represented as a separate standalone component as before.
In these platforms, inter-worker communication is often realized implicitly through shared-memory reads and writes within the same address space, while atomic operations enforce fine-grained synchronization.
As a result, communication cost appears as memory-system pressure rather than explicit message passing.
For traversal-based graph analytics, which feature irregular access patterns and low arithmetic intensity, performance is often bounded by data movement through the memory hierarchy rather than by arithmetic operations.
Moreover, because frontier sizes and degree distributions vary across natural graphs, achievable performance also depends on whether sufficient parallel work can be sustained throughout execution \cite{gonzalez2012powergraph}.

% Having inspired by the BSP model and architectural advances of multi-core CPUs and GPUs, 
Motivated by the limitations of BSP and by architectural advances in multi-core CPUs and GPUs, we introduce a profiling-based methodology to evaluate execution efficiency along two complementary dimensions: data movement effectiveness and parallelism exploitation. For each dimension, we define metrics that expose bottlenecks and enable a CPU--GPU comparison of memory and compute utilization. We quantify data movement using CPU cache behavior (L3 miss rate and miss frequency, reflecting LLC utilization) and GPU effective memory bandwidth (attainable device-memory throughput). We quantify parallelism exploitation using instruction execution rate on both CPU and GPU, indicating how effectively the hardware issues and retires instructions given the available parallel work. Together, these metrics support resource-utilization and platform-suitability assessment by indicating whether performance is limited by data movement or insufficient parallel work.

\subsection{An Executive Summary of Our Work}\label{subsec:intro-exec-summary}
In this work, we study four representative traversal-based graph algorithms: BFS, SSSP, MSSP, and the blossom algorithm for maximum matching.
Using our profiling-based methodology, we conduct a comprehensive performance- and measurement-driven study to understand their execution behavior on multi-core CPUs and GPUs, extract key performance insights, and assess their architecture suitability across the two platforms.
Our contributions are summarized as follows:
\begin{itemize}[leftmargin=*, topsep=0pt, itemsep=0pt, parsep=0pt]
  \item 
  \textcolor{blue}{We add a selective-reset mechanism for alternating trees to the existing X-Blossom implementation}, resulting in an improved CPU version, X-Blossom-Pro, which accelerates the original X-Blossom by up to $3.26\times$.
  \item 
  We develop the first GPU-based implementation of a blossom-style maximum matching algorithm, X-Blossom++, \textcolor{blue}{by porting the dynamic data structures such as path table and checking set to GPU-friendly buffers while using fine-grained load balancing strategy to maximize parallelism exploitation. }
  As a result, X-Blossom++ achieves speedups of up to $39.09\times$ over X-Blossom-Pro.
  \item We introduce architecture-aware profiling metrics to characterize traversal-based graph workloads and assess CPU vs.~GPU suitability.
  \item \phantomsection\label{item:takeaway-architecture-suitability}\textcolor{blue}{Our analysis yields the architectural-suitability takeaway: workloads with high parallelism in both computation and data movement benefit more from GPUs, whereas workloads with limited parallelism and small, high-locality working sets are better suited to multi-core CPUs.}
  \item The datasets, program implementations on both CPUs and GPUs, and scripts to reproduce all experiments are open sourced for public use and continued development at \href{https://github.com/victorliu-sq/GACGE-SIGMETRICS27}{this link}.
\end{itemize}
